\subsection*{Level distribution is a value-distribution and packet input (NS-29)}

The depth is $D_a=\| (\beta_a)_-\|_\infty=-\inf\beta_a$ when
$\inf\beta_a<0$, not the signed infimum itself. Fix $a>0$, an integer
head size $N\ge1$, and $X=2\pi N/L$, where $L=2a$. The head cutoff
specifies the observation interval $[0,X]$; it does not truncate the
complete symbol, the Fourier transform, or the form. In any cofinal
assertion below $N=N(a)$ must be specified as part of the family.
In particular the global depth $D_a$ must not be replaced silently by
$D_{a,X}:=\max(0,-\min_{[0,X]}\beta_a)\le D_a$.

\begin{proposition}[Separation of level measure and packet coverage]
\label{prop:v147-level-separation}
Fix $D=D_a>0$ and $0<\theta\le1$. Put
\[
 F_{a,X}(t)=|\{\xi\in[0,X]:\beta_a(\xi)<-t\}|,
 \qquad E(t_1,t_2)=\{\xi\in[0,X]:-t_2<\beta_a(\xi)<-t_1\}.
\]
The following is an explicit sufficient level-measure hypothesis:
\begin{equation}\label{eq:v147-level-increments}
 F_{a,X}(t_1)-F_{a,X}(t_2)
 =|E(t_1,t_2)|\ge\kappa X\frac{t_2-t_1}{D}
 \quad(0\le t_1<t_2\le\theta D),\qquad\kappa>0.
\end{equation}
Suppose also that a source-admissible even unit vector
$g\in\mathcal D(q_a)$, a measurable $S\subset[0,X]$, and constants
$b,\alpha>0$, $Q\ge0$ satisfy
\begin{align}
 |\mathcal F_ag(\xi)|^2&\ge b/X\quad\hbox{a.e. on }S,
 &q_a[g]&\le Q, \notag\\
 |S\cap E(t_1,t_2)|&\ge\alpha|E(t_1,t_2)|
 &&(0\le t_1<t_2\le\theta D).
 \label{eq:v147-packet-coverage}
\end{align}
Then, for every Borel $B\subset[-\theta D,0]$,
\begin{equation}\label{eq:v147-weighted-density}
 \mu_g(B)\ge\frac{2b\alpha\kappa}{D}|B|.
\end{equation}
For $\theta=1$ this is precisely the hypothesis of
Proposition~\ref{prop:v146-level-complexity}, with $c=2b\alpha\kappa$.
For $0<\theta<1$ it is a weaker, sufficient replacement: every minorant
with at most $K$ distinct values satisfies
\begin{equation}\label{eq:v147-partial-level-loss}
 \eta(m)\ge\left[\frac{b\alpha\kappa\theta^2D}{K}-Q\right]_+.
\end{equation}
Thus uniform positive $b,\alpha,\kappa,\theta$ and bounded $Q$ on a
cofinal family imply the same linear-in-$D_a$ level-count obstruction.
These are additional hypotheses, not conclusions about the arithmetic
symbol. Taking $S=[0,X]$ permits $\alpha=1$. Merely requiring
$|S|\ge\alpha X$, or an upper bound on Fourier density, does not
replace \eqref{eq:v147-packet-coverage}.
\end{proposition}
\begin{proof}
The exact $\beta_a$ is nonconstant and real analytic on the real axis,
and tends to $+\infty$ at infinity. Its level sets in $[0,X]$ have
Lebesgue measure zero, so the equality in
\eqref{eq:v147-level-increments} is valid, including at the endpoints.
On each band, integration of the Fourier lower bound gives at least
$b\alpha\kappa(t_2-t_1)/D$. Evenness doubles this contribution on the
reflected interval; all omitted frequency mass is nonnegative.
Domination on intervals extends to Borel sets (equivalently, apply the
Radon--Nikodym derivative of the level pushforward). This proves
\eqref{eq:v147-weighted-density}.

For the last bound partition $[-\theta D,0]$ at the minorant's values
in its interior. There are at most $K$ gaps: at least one value is
$-D\le-\theta D$, since the minorant is at least $-D$, has finitely
many values, and is below a symbol with infimum $-D$. On each gap the
minorant is at most its lower endpoint; if that endpoint is the
inserted $-\theta D$, this remains an upper bound. For gap lengths
$\Delta_j$ the loss is therefore at least
\[
 \frac{b\alpha\kappa}{D}\sum_j\Delta_j^2
 \ge\frac{b\alpha\kappa\theta^2D}{K}.
\]
Subtract $q_a[g]\le Q$ using \eqref{eq:v146-gap-packet}.
Unconditional $D_a\to\infty$ proves the conditional cofinal assertion.
\end{proof}

A cumulative estimate such as
$F_{a,X}(t)\ge\kappa X(1-t/D)$ is not enough. For example, the
nonarithmetic affine symbol $\beta(\xi)=-D/2-D\xi/(2X)$ has
$F(t)=X$ for $0<t\le D/2$ and $F(t)=2X(1-t/D)$ for $D/2<t<D$.
It satisfies that cumulative estimate with $\kappa=1$, but has no
mass in any level interval contained in $(-D/2,0)$. This is a
counterexample to the measure-theoretic inference, not to the
arithmetic symbol. Subtracting two lower bounds is not a lower bound
for their difference. Likewise two sets each having positive measure
need not intersect in every level band.

There is also an exact local description. At almost every regular level
$y$, the unweighted and weighted densities on $[0,X]$ are
\begin{equation}\label{eq:v147-coarea}
 h_{a,X}(y)=\sum_{\substack{0<\xi<X\\\beta_a(\xi)=y}}
          \frac1{|\beta_a'(\xi)|},\qquad
 h_{g,X}(y)=\sum_{\substack{0<\xi<X\\\beta_a(\xi)=y}}
          \frac{|\mathcal F_ag(\xi)|^2}{|\beta_a'(\xi)|}.
\end{equation}
This follows by splitting at the finitely many critical points and
changing variables on the monotone pieces. The increment condition is
equivalent to $h_{a,X}(y)\ge\kappa X/D$ almost everywhere on the
specified level interval. It asks for enough crossings with controlled
slopes, not just zeros of $\zeta$ or one deepest trough. Necessarily
$D_{a,X}\ge\theta D_a$; the finite head need not see the global depth.

This separation is sufficient, not necessary for the original weighted
hypothesis. The single-interval geometry following
Proposition~\ref{prop:v146-level-complexity} could supply the weighted
level density on a frequency interval of length $O(1/a)$ with packet
density of order $a$, without a positive fraction of the whole $[0,X]$.
There is also no assertion uniform over arbitrary choices of $N$:
for fixed $a$ the negative set of $\beta_a$ is bounded, so
$F_{a,X}(0)/X\to0$ as $X\to\infty$. Choosing excessively large
heads can therefore defeat \eqref{eq:v147-level-increments} without
affecting the complete weighted condition. Prescribing the cofinal
range is part of ZLD, not a harmless choice made after the estimate.

\paragraph{Exact reduction to a zero statement.}
Use all nontrivial zeros with multiplicity and the symmetric-height
Perron convention of Proposition~\ref{prop:v146-beta-explicit}. Define
\begin{align}
 Z_a(\xi)&=2\Re\lim_{T\to\infty}
       \sum_{|\Im\rho|<T}\frac{x^{\rho-1/2+i\xi}}{\rho-1/2+i\xi},
 &x&=e^{2a}, \notag\\
 C_a(\xi)&=e_x\cos(2a\xi)
       -2\Re\sum_{n\ge1}\frac{x^{-2n-1/2+i\xi}}{2n+1/2-i\xi},
 &\mathcal Z_a(\xi)&=Z_a(\xi)+C_a(\xi)=\beta_a(\xi).
 \label{eq:v147-corrected-zero-field}
\end{align}
Real continuous limits are used at critical-line zeros. No zero-location
hypothesis enters this identity. In particular the endpoint is retained
for prime-power $x$. Uniformly in real $\xi$,
\begin{equation}\label{eq:v147-correction-size}
 |C_a(\xi)|\le\varepsilon_a:=
 e_x+\frac{4x^{-5/2}}{5(1-x^{-2})}\longrightarrow0.
\end{equation}
Indeed $|2n+1/2-i\xi|\ge2n+1/2\ge5/2$, and
$e_x\le(\log x)/\sqrt x$. Consequently
\[
 \{Z_a<-t-\varepsilon_a\}\subset\{\beta_a<-t\}
 \subset\{Z_a<-t+\varepsilon_a\}.
\]
For a band one must instead shrink both ends:
\[
 \{-t_2+\varepsilon_a<Z_a<-t_1-\varepsilon_a\}
 \subset E(t_1,t_2)
 \quad\hbox{within }[0,X].
\]
The inner band can be empty when $t_2-t_1\le2\varepsilon_a$.
Thus even a small correction cannot simply be discarded in a demand
about \emph{every} level interval. A finite zero prefix also needs the
Perron remainder $-2\Re R_{x,T}$; \eqref{eq:v147-correction-size} is
not an estimate for that remainder.

\begin{assumption}[Corrected zero level distribution, ZLD]
\label{ass:v147-zld}
Fix a cofinal set of windows, a prescribed head-size function $N(a)$,
and $0<\theta\le1$. There is a constant $\kappa>0$ independent of
$a$ such that, eventually, for every $0\le r_1<r_2\le\theta$,
\begin{equation}\label{eq:v147-zld}
 \frac1{X_a}\left|\left\{0\le\xi\le X_a:
  -r_2D_a<\mathcal Z_a(\xi)<-r_1D_a\right\}\right|
 \ge\kappa(r_2-r_1),\qquad X_a=2\pi N(a)/L.
\end{equation}
This is a named \emph{open input}, not an established conjecture from
the literature. The independent packet input is
\eqref{eq:v147-packet-coverage} with $b,\alpha$ bounded below and
$Q$ bounded above on that same family. Neither input is asserted here.
\end{assumption}
Substituting \eqref{eq:v147-corrected-zero-field} proves that ZLD is
exactly \eqref{eq:v147-level-increments}. No information is gained by
changing from the prime expression to the zero expression. Only ZLD
with $\theta=1$ addresses the full original interval $[-D_a,0]$;
fixed $\theta<1$ suffices for the modified bound
\eqref{eq:v147-partial-level-loss}.

\paragraph{What the zero literature supplies, and the missing reduction.}
The following assessment concerns the stated theorems, not a proof
that no future consequence of zero theory can yield ZLD.

Riemann--von Mangoldt gives
$N(T)=\frac{T}{2\pi}\log\frac{T}{2\pi e}+O(\log T)$; explicit
versions are given by Hasanalizade--Shen--Wong
\cite{HasanalizadeShenWong2022}. It counts ordinates, without a signed
value-distribution conclusion for \eqref{eq:v147-corrected-zero-field}.
At large height $T$, an interval of width $O(1/a)$ has main count
$O((\log T)/a)$, while differencing the displayed error leaves
$O(\log T)$. In particular this count estimate alone does not isolate
or control cancellation at every window-scale side lobe.

Zero-density theorems bound the number of zeros off a specified
vertical line. For example Ingham's bound is
$N(\sigma,T)\ll T^{3(1-\sigma)/(2-\sigma)}\log^5T$; see
\cite{Ingham1940,ChourasiyaSimonic2025}. This is an upper count,
not a lower signed band measure. In the present sum an off-line zero
also has the factor $x^{\Re\rho-1/2}$, which may be large. Replacing
all those factors by one uses RH. Removing this issue by assuming RH
still leaves phase cancellation and the normalization by the extreme
value $D_a$; it does not supply a proof of ZLD.

Landau--Gonek estimates control the linear sums
$\sum_{0<\Im\rho\le T}x^\rho$, including an arithmetic main term
at prime powers; see \cite[Introduction, (1.1)]{Aryan2019}.
They do not state a signed distribution theorem for a moving resolvent.
Even on RH, with $A_x(v)=\sum_{0<\gamma\le v}x^{i\gamma}$ and
$\xi<U<V$, partial summation reads
\[
 \sum_{U<\gamma\le V}\frac{x^{i\gamma}}{\gamma-\xi}
 =\frac{A_x(V)}{V-\xi}-\frac{A_x(U)}{U-\xi}
       +\int_U^V\frac{A_x(v)}{(v-\xi)^2}\,dv.
\]
An upper error bound here loses control as the interval approaches
$\xi$ and gives no lower signed quantile. The real, paired expression
has removable singularities; the problem is the loss of its cancellation
in such absolute estimates, not a pole in $\beta_a$. Controlling a
first or second moment likewise does not give a lower density on
\emph{every} interval of values, particularly relative to a supremum.

Montgomery's theorem and pair-correlation conjecture concern averages
of differences of ordinates \cite{Montgomery1973}; the unconditional
complex-zero extension also remains a correlation average
\cite[Theorem 1]{BaluyotEtAl2023}. ZLD instead concerns values of a
phase-sensitive complete sum, with $x=e^{2a}$ varying, on a prescribed
frequency interval and normalized by its global negative depth. No
reduction from these correlation statements to ZLD, or conversely, is
proved here or supplied by the cited results. The same applies to the
classical density hypothesis
$N(\sigma,T)\ll_\epsilon T^{2(1-\sigma)+\epsilon}$ (and to
Lindel\"of's hypothesis): these control counts or upper growth, not the
required signed level distribution. ZLD is therefore a separate input
with \emph{unclassified logical strength}; it is not justified to call
it weaker than RH, equivalent to RH or pair correlation, or logically
independent of them. Its conjunction with the packet input obstructs
a minorant method; it is not a positivity assertion.

\begin{proposition}[What depth growth actually says about measure]
\label{prop:v147-depth-measure}
Unconditionally, at a fixed window define
\[
 M_a=\tfrac12\psi'(5/4)
       +2\sum_{1<n<x}\frac{\Lambda(n)\log n}{\sqrt n}
       +2\int_0^L t e^{t/2}\,dt.
\]
Then $|\beta_a'|\le M_a$. If $d=D_{a,X}>t\ge0$, then
\begin{equation}\label{eq:v147-depth-width}
 F_{a,X}(t)\ge\min\{X,(d-t)/M_a\}.
\end{equation}
If $\beta_a([0,X])$ contains $[-\theta D_a,0]$, then
\eqref{eq:v147-level-increments} holds at this window with
$\kappa=D_a/(M_aX)$. No uniform positive lower bound for this
$\kappa$ follows from $D_a\to\infty$.

Under the additional RH hypothesis of
Proposition~\ref{prop:v137-conditional-mass}, let
$w_a(\xi)=w(2a(\xi-\gamma_0))$. Eventually, for $0\le t<D_a$,
\begin{equation}\label{eq:v147-probe-measure}
 |\{\beta_a<-t\}\cap\operatorname{supp}w_a|
 \ge\frac{[1/(6\pi)-\pi t/a]_+}{D_a-t}.
\end{equation}
This is a conditional, possibly zero, cumulative measure bound. It
is not ZLD or a physical packet estimate.
\end{proposition}
\begin{proof}
Differentiate the finite prime sum and the continuum integral in
\eqref{eq:v130-r-symbol}. The trigamma series implies
$|\psi'(5/4+i\xi/2)|\le\psi'(5/4)$, proving the derivative bound.
From a minimizer on $[0,X]$, the Lipschitz estimate supplies a
one-sided interval of length at least the right side of
\eqref{eq:v147-depth-width}; strict endpoints have zero measure.
If the range contains the indicated levels, almost every such level
has a regular preimage. Equation~\eqref{eq:v147-coarea} then gives
$h_{a,X}\ge1/M_a$. This proves the fixed-window increment estimate.

For the RH statement, $0\le w_a\le1$, $\int w_a=\pi/a$ and,
eventually, $J_a=\int\beta_aw_a\le-1/(6\pi)$. Splitting the
integral at $\beta_a=-t$ gives
\[
 J_a\ge-\frac{\pi t}{a}
        -(D_a-t)\int_{\{\beta_a<-t\}}w_a.
\]
Rearranging and using $w_a\le1$ proves the claim.
\end{proof}
At $t=\theta D_a$, the numerator in
\eqref{eq:v147-probe-measure} need not be positive: the argument has
no upper control on $D_a/a$. Even a positive value is not a bound on
all band differences normalized by $X_a$. The unconditional v1.37
argument is a contradiction: a bounded cofinal scalar floor would
imply RH and then contradict the conditional probe. It proves global
depth divergence, not a uniform sublevel measure, a head-range depth,
or the Fourier mass of a realizable packet. Continuity supplies local
width only with its window-dependent slope cost, as made explicit above.

\begin{proposition}[Broad Fourier coverage does not bound the energy]
\label{prop:v147-broad-packet}
Fix a nonnegative, even $h\in C_c^\infty(-1/2,1/2)$ of $L^2$ norm
one. For $X\to\infty$ and $a>1/(2X)$ put
$f_X(t)=\sqrt X h(Xt)$, a unit even physical window vector. For a
constant $b_h>0$ depending only on $h$,
\begin{equation}\label{eq:v147-broad-packet}
 |\mathcal F_af_X(\xi)|^2\ge b_h/X\quad(0\le\xi\le X),
 \qquad q_a[f_X]=\log X+C_h+o(1).
\end{equation}
The energy asymptotic is uniform in these windows once
$1/X<\log2$. It gives no uniform $Q$ for this construction.
These vectors are not asserted source-admissible.
\end{proposition}
\begin{proof}
The unitary transform is $X^{-1/2}\widehat h(\xi/X)$. On $|v|\le1$,
$\widehat h(v)\ge(2\pi)^{-1/2}\cos(1/2)\int h>0$, proving the
lower bound. The autocorrelation of $f_X$ is supported in
$[-1/X,1/X]$. Every prime lag $\log n\ge\log2$ therefore contributes
zero exactly; the continuum term contributes $O(1/X)$, since the
autocorrelation has absolute value at most one. Finally
\[
 \int_{\mathbb R}A(Xv)|\widehat h(v)|^2\,dv
 =\log X-\log(2\pi)
       +\int_{\mathbb R}\log|v|\,|\widehat h(v)|^2\,dv+o(1).
\]
This follows from the digamma asymptotic; split $|v|<1/X$, whose
logarithmic contribution is $O((\log X)/X)$, and use the rapid decay
of $\widehat h$ on the remainder. The displayed integral defining
$C_h$ is finite. No prime cancellation or RH is used.
\end{proof}

Source projection has a pointwise cost as well as an energy cost. For
any unit even $f\in\mathcal D(q_a)$, with the actual source
$u_a\in\mathcal D(\mathsf W_{e^a})$, write $r=|\langle f,u_a\rangle|<1$ and
$g=P_af/\sqrt{1-r^2}$. Then
\begin{equation}\label{eq:v147-projected-density}
 |\mathcal F_ag(\xi)|\ge
 \frac{[|\mathcal F_af(\xi)|-r|\mathcal F_au_a(\xi)|]_+}
      {\sqrt{1-r^2}}.
\end{equation}
In particular a raw $b/X$ lower bound persists with $b/(4X)$ if
$r\|\mathcal F_au_a\|_\infty\le\sqrt{b/X}/2$. Small ordinary
source overlap alone is not this scale-dependent inequality.
The exact form expansion also gives
\[
 \left|(1-r^2)q_a[g]-q_a[f]\right|
 \le(2r+r^2)\|\mathsf W_{e^a}u_a\|.
\]
Thus for the broad pulses, if $r\to0$ and the source residual is
bounded, projection still leaves $q_a[g]/\log X\to1$.
It cannot repair their energy cost. The packet of a single trough in
Lemma~\ref{lem:v146-packet-concentration}, conversely, supplies local
Fourier mass but not coverage of every level across $[0,X]$ or a
uniform original energy. At one fixed window every nonzero physical
packet has an analytic transform and hence can be bounded below off
small neighborhoods of its finitely many zeros in $[0,X]$; the bound
and the energy may depend on the window, and those neighborhoods
can meet precisely the needed level bands. This does not give
\eqref{eq:v147-packet-coverage}.

\paragraph{Finding and scope.}
The outcome of NS-29 is a reduction to the named open input ZLD
\emph{together with} bounded-energy packet coverage. The existing
unconditional zero theorems assessed above do not complete that reduction
into a proved bound. Nor can one identify packet spread alone as the
remaining obstacle: the uniform arithmetic band measure (a) is itself
unproved, and even granting it leaves the joint coverage and energy
requirement (b). The modification from cumulative sublevel size to
\emph{increments}, and from an arbitrary positive fraction to
\emph{coverage of each level band}, is essential.

The new fixed-window implications require no RH assumption. RH is used
explicitly in \eqref{eq:v147-probe-measure} and when interpreting the
critical-line exponential sums in the literature discussion. The two cofinal
inputs remain open. The argument and factors of two apply to odd
packets as well because their Fourier modulus is even; the even-source
constraint is then absent. NS-22 and the concurrent NS-28 measurements
are diagnostics, not substitutes for these cofinal hypotheses. No new
window, head computation, tail metric or zero enclosure is made. This
is a gap in a proposed minorant obstruction, not failure of the true
weighted concentration criterion or of the physical object. The route
remains blocked, not closed. G2 and RH remain open.

